\documentclass[11pt,a4paper]{article}

\usepackage[T1]{fontenc}
\usepackage[utf8]{inputenc}
\usepackage[czech]{babel}
\usepackage{lmodern}
\usepackage{microtype}
\usepackage[margin=3cm]{geometry}
\usepackage{amsmath,amssymb}
\usepackage{enumitem}
\usepackage{listings}
\usepackage{xcolor}
\usepackage[hidelinks]{hyperref}

% --- prostředí pro invarianty -------------------------------------------
% Stabilní číslování INV-01, INV-02, ... Číslo se nesmí měnit ani při
% přeuspořádání textu -- invarianty jsou citovatelné zvenčí.
\newcounter{invariant}
\renewcommand{\theinvariant}{INV-\ifnum\value{invariant}<10 0\fi\arabic{invariant}}
\newenvironment{invariant}[1]
  {\refstepcounter{invariant}\par\medskip\noindent
   \textbf{\theinvariant\quad #1}\par\nopagebreak\smallskip\itshape}
  {\par\medskip}

\lstset{
  basicstyle=\ttfamily\small,
  breaklines=true,
  frame=leftline,
  framesep=6pt,
  xleftmargin=8pt,
  showstringspaces=false,
}

\title{Dodávka on-premises produktu\\\large Model divergence stavu a mechanismy jeho zúžení}
\author{Jiří Slachta}
\date{Pracovní verze --- \today}

\begin{document}
\maketitle

\begin{abstract}
\noindent
Text popisuje dodávku software provozovaného v prostředí zákazníka jako
problém řízení divergence: rozdílu mezi skutečným stavem instance a
představou dodavatele o tomto stavu. Ukazuje, že náklad této formy dodávky
neroste s rychlostí vývoje, ale s počtem živých kombinací, které musí
dodavatel obsloužit, a že většina běžně doporučovaných opatření
(pomalá kadence, minimalizace počtu vydání) tento počet ve skutečnosti
zvyšuje. Jádrem je model cestové závislosti stavu databáze a z něj
odvozená sada invariantů a mechanismů.
\end{abstract}

\tableofcontents
\bigskip

% ========================================================================
\section{Rámec}
\label{sec:ramec}

\emph{[Stub. Obsah: on-prem jako provozní režim, nikoli distribuční model.
Rozlišení mezi tím, co skutečně zaniklo (ruční instalace bez pipeline),
a tím, co přetrvává (regulace, OT, ekonomika stabilní zátěže, inference
nad citlivými daty). Nákladový driver: počet živých kombinací, nikoli
rychlost vývoje. Vyvrácení intuice, že pomalá kadence zužuje pole verzí.]}

% ========================================================================
\section{Model}
\label{sec:model}

\subsection{Dva stavy}

Instance provozovaná u zákazníka má v každém okamžiku dva stavy, které
je nutné držet odděleně.

\textbf{Skutečný stav} $S$ je to, co je fyzicky v databázi a na disku:
struktura schématu, obsah dat, přítomnost či nepřítomnost objektů, které
tam kdo kdy vytvořil.

\textbf{Evidovaný stav} $E$ je to, co si o instanci myslí dodavatel a
co si o sobě myslí instalátor: číslo verze, obsah migračního ledgeru,
záznamy v systémech dodavatele.

V ideálním případě platí $S \equiv E$. V praxi vzniká divergence
$D = S \setminus E$ --- množina skutečností, které v instanci platí,
ale nikde nejsou evidované. Tato práce vychází z jediné teze:

\begin{quote}
\textbf{Selhání dodávky on-premises produktu nevzniká z chyb v kódu.
Vzniká tehdy, když se rozhoduje podle $E$, zatímco operace působí na $S$.}
\end{quote}

Migrace napsaná proti evidovanému stavu spadne, protože je aplikována na
stav skutečný. Rozhodnutí o vydání učiněné podle evidence o verzích v poli
je chybné, protože pole vypadá jinak. Diagnóza vzdálené instance je chybná,
protože se opírá o číslo verze, které nepopisuje to, co je uvnitř.

Všechny mechanismy popsané v sekci~\ref{sec:mechanismy} mají jediný účel:
zmenšit $D$, případně učinit ji pozorovatelnou dřív, než se o ni někdo
opře.

\subsection{Verze schématu není stav systému}
\label{sub:pathdep}

Většina migračních nástrojů drží jediné číslo --- kde jsem. Tím implicitně
tvrdí, že stav je funkcí cíle:
\[
S = f(v)
\]
kde $v$ je verze schématu. Toto tvrzení je nepravdivé. Stav je funkcí
\emph{cesty}:
\[
S = f\bigl(S_0,\ \langle m_1, m_2, \dots, m_n \rangle\bigr)
\]
kde $S_0$ je výchozí stav a $\langle m_i \rangle$ posloupnost skutečně
aplikovaných transformací v pořadí, ve kterém proběhly.

Dvě instance na téže verzi tedy nemusí být v témže stavu. Instance, která
prošla $v_3 \to v_4 \to \dots \to v_7$, a instance, která přišla jinou
cestou nebo některé kroky přeskočila, mají shodné $v$ a odlišné $S$.

Praktický důsledek: \emph{prostor, který musí dodavatel testovat, není
množina verzí, ale množina cest.} Ta roste podstatně rychleji a na rozdíl
od množiny verzí není nikde evidovaná.

\subsection{Konvergentní a divergentní transformace}

Ne každá migrace přispívá k divergenci. Rozhodující je, zda je její
výsledek závislý na vstupním stavu.

\textbf{Konvergentní transformace} má výsledek nezávislý na $S$ i na cestě.
Typicky čistě strukturální operace:

\begin{lstlisting}
ALTER TABLE WorkOrder ADD COLUMN ClosedAt datetime2 NULL;
\end{lstlisting}

\noindent
Ať byla instance v jakémkoli stavu, po této operaci má sloupec. Cesty
konvergují.

\textbf{Divergentní transformace} má výsledek, který je funkcí vstupního
stavu:

\begin{lstlisting}
UPDATE WorkOrder
   SET ClosedAt = (SELECT MAX(ChangedAt) FROM WorkOrderLog l
                    WHERE l.WorkOrderId = WorkOrder.Id
                      AND l.NewState = 'Closed');
\end{lstlisting}

\noindent
Výsledek závisí na tom, co bylo v \texttt{WorkOrderLog} v okamžiku běhu.
Dvě instance se stejným schématem a odlišnou historií skončí s odlišnými
daty. Cesty divergují.

Hranice tedy nevede mezi \uv{malými} a \uv{velkými} migracemi, ale mezi
těmi, které se dotýkají pouze struktury, a těmi, které se dotýkají dat.
Každá data-touching migrace je potenciální zdroj divergence a musí být
posuzována jinou optikou.

\subsection{Podmínky bezpečnosti divergentní transformace}

Data-touching migraci lze učinit bezpečnou, pokud splňuje tři podmínky
současně.

\begin{description}[leftmargin=2em,style=nextline]
\item[Determinismus]
Nad týmž vstupem dává týž výstup. Zejména nezávisí na čase běhu,
na lokálním nastavení instance ani na pořadí, ve kterém databáze
vrátila řádky.

\item[Idempotence]
Opakované spuštění nezmění výsledek. Druhý běh se dotkne nula řádků.

\item[Čistota vzhledem ke vstupnímu stavu]
Výsledek je odvoditelný výhradně z dat, která jsou v okamžiku běhu
v databázi. Nikoli z toho, jaká byla předchozí verze aplikace nebo
který patch instance minula.
\end{description}

Pokud transformaci nelze takto napsat, není to migrace. Je to předpoklad,
a předpoklad patří do preflightu jako ověřovaný invariant, nikoli do
mutační fáze.

\subsection{Zdroje nedeterminismu, které projdou code review}
\label{sub:nedeterminismus}

Nedeterminismus bývá v migracích neúmyslný a vizuálně nenápadný.
Opakující se zdroje:

\begin{itemize}[leftmargin=1.4em]
\item \textbf{Aktuální čas.} \texttt{GETDATE()} v backfillu. Migrace
      spuštěná o rok později vyrobí jiná data než migrace spuštěná
      v den vydání.
\item \textbf{Konfigurace instance.} Locale, časová zóna, collation,
      výchozí měna, výchozí jednotky. Tytéž řádky, jiný výsledek.
\item \textbf{Objem dat v okamžiku běhu.} Zejména u transformací
      pracujících po dávkách, kde se počet dávek propisuje do výsledku.
\item \textbf{Nedeterministické pořadí.} \texttt{TOP} nebo \texttt{LIMIT}
      bez úplného \texttt{ORDER BY}. Databáze nemá povinnost vrátit
      dvakrát totéž.
\item \textbf{Volání aplikační logiky.} Migrace, která pro výpočet
      použije kód aplikace, je svázaná s verzí toho kódu. Ta se mění,
      migrace zůstává.
\end{itemize}

Společný jmenovatel: všechny tyto zdroje jsou neviditelné v okamžiku,
kdy se migrace píše, protože autor ji spouští jednou, ve svém prostředí,
ihned.

\subsection{Dvě cesty k témuž schématu}
\label{sub:freshvsupgrade}

Zvláštní případ cestové závislosti, který se projeví až po letech.

Nová instalace obvykle vzniká z baseline skriptu. Existující instance
došla řetězem migrací. Obě jsou nominálně na téže verzi a obě jsou
funkční. Liší se v detailech, které nikoho nezajímají, dokud na ně něco
nesáhne:

\begin{itemize}[leftmargin=1.4em]
\item jména omezení --- explicitní \texttt{PK\_WorkOrder} proti
      autogenerovanému \texttt{PK\_\_WorkOrde\_\_3213E83F\dots}
\item pořadí sloupců v tabulce
\item existence indexů, které vznikly jako vedlejší efekt starší migrace
      a v baseline nejsou
\item \texttt{IDENTITY} seedy a mezery v číselných řadách
\item nullability a collation na úrovni sloupce, kde se chování
      \texttt{ALTER} liší od \texttt{CREATE}
\end{itemize}

Pak přijde migrace obsahující
\texttt{DROP CONSTRAINT [FK\_WorkOrder\_Asset]} a selže u té poloviny pole,
která má omezení pojmenované jinak.

Tento případ je uveden zvlášť, protože je na rozdíl od ostatních
\emph{levně a úplně automatizovatelně detekovatelný} --- viz
sekce~\ref{sec:mechanismy}.

\subsection{Latence detekce jako hlavní nákladový faktor}

Divergence se neprojeví v okamžiku vzniku. Projeví se při první operaci,
která se o porušený předpoklad opře. Mezi těmito dvěma okamžiky uplyne
libovolně dlouhá doba.

To má tři důsledky, které určují ekonomiku celé dodávky:

\begin{enumerate}[leftmargin=1.6em]
\item \textbf{Náklad opravy roste s latencí.} Divergence odhalená před
      upgradem je konfigurační problém. Táž divergence odhalená během
      upgradu je havárie s odstávkou. Odhalená tři měsíce po upgradu je
      forenzní úloha nad poškozenými daty.
\item \textbf{Příčina se s latencí ztrácí.} Po dostatečné době není
      zjistitelné, která migrace stav porušila, protože chybí záznam
      o tom, co bylo předtím.
\item \textbf{Divergence se skládá.} Neopravená divergence se stává
      vstupním stavem pro další transformace, které s ní nepočítají.
\end{enumerate}

Cílem tedy není nulová divergence --- ta je nedosažitelná. Cílem je
\textbf{krátká latence detekce}: aby se rozpor mezi $S$ a $E$ projevil
v okamžiku, kdy je jeho oprava levná, a nikoli v okamžiku, kdy je drahá.

\subsection{Co z modelu plyne pro měřitelnost}

Aby byla latence krátká, musí být divergence pozorovatelná. Z předchozích
odstavců plynou tři veličiny, které musí být zaznamenatelné a přenositelné:

\begin{description}[leftmargin=2em,style=nextline]
\item[Cesta, nikoli verze]
Ledger aplikovaných transformací: co, kdy, jakou verzí nástroje, s jakým
výsledkem, s jakým otiskem skriptu v okamžiku aplikace. Otisk je nutný,
protože skript opravený po vydání dělí pole na instance před opravou
a po ní.

\item[Otisk dat, nikoli jen schématu]
Kontrolní součty nad klíčovými tabulkami, distribuce hodnot výčtových
sloupců, počty porušení jednotlivých invariantů, seznam objektů v databázi,
které nepocházejí od dodavatele.

\item[Výsledek nezávislého ověření]
Tvrzení o platnosti invariantů získané dotazem nad daty, nikoli odvozené
z evidence o tom, že migrace proběhla.
\end{description}

Poslední bod je podstatou celého modelu. Evidence o provedení operace
není důkazem o jejím výsledku.

% ========================================================================
\section{Invarianty}
\label{sec:invarianty}

\emph{[Stub. Číslovaná, citovatelná sada. Každý invariant: jedna věta
tvrzení, odstavec zdůvodnění odkazující do sekce~\ref{sec:model},
poznámka o tom, jak se ověřuje. Číslo je stabilní a nemění se při
přeuspořádání textu.]}

\begin{invariant}{Schéma se v rámci patch řady nemění.}
\emph{[Zdůvodnění, výjimka pro aktivně destruktivní stav, způsob ověření.]}
\end{invariant}

\begin{invariant}{Data-touching migrace je deterministická, idempotentní a čistá vzhledem ke vstupnímu stavu.}
\emph{[Odkaz na \ref{sub:nedeterminismus}. Co dělat, když to nejde napsat.]}
\end{invariant}

\begin{invariant}{Každá data-touching migrace ověřuje vlastní úplnost.}
\emph{[Post-condition tvaru \uv{zbývá nula}, nikoli \uv{update proběhl}.]}
\end{invariant}

\begin{invariant}{Baseline a plný řetěz migrací vedou k identickému schématu.}
\emph{[Odkaz na \ref{sub:freshvsupgrade}. Ověřuje se schema diffem v CI.]}
\end{invariant}

\begin{invariant}{Poslední patch v řadě je jediným podporovaným vstupem do dalšího majoru.}
\emph{[Zúžení vstupní brány. Vynucuje instalátor, nikoli dokumentace.]}
\end{invariant}

\begin{invariant}{Destruktivní změna se nikdy neprovádí v jednom vydání.}
\emph{[Expand/contract napříč verzemi. Bod bez návratu neexistuje.]}
\end{invariant}

% ========================================================================
\section{Mechanismy}
\label{sec:mechanismy}

\emph{[Stub. Každý mechanismus explicitně uvádí, který invariant vynucuje
a jakou divergenci zkracuje. Kandidáti:]}

\begin{itemize}[leftmargin=1.4em]
\item \emph{Rozdělení migrace na preflight a mutate.} Read-only ověření
      předpokladů spustitelné kdykoli, i měsíce před vydáním.
\item \emph{Expand/contract napříč vydáními.}
\item \emph{Opravy jako idempotentní repeatable skripty} namísto položek
      v lineárním řetězu.
\item \emph{Konvergenční uzávěrka} na konci patch řady.
\item \emph{Schema diff v CI} přes všechny reálné cesty.
\item \emph{Profil instance jako telemetrie} --- tvary dat, nikoli data.
\item \emph{Kohorty (rings)} pro postupné uvolňování.
\end{itemize}

% ========================================================================
\section{Rozhodnutí, která zůstávají úsudkem}
\label{sec:usudek}

\emph{[Stub. Nikoli pravidla, ale vyjmenované veličiny, které se váží.
Cílem kapitoly je odlišit úsudek nad známými veličinami od intuice.
Kandidáti: hotfix versus odklad do nejbližšího vydání; prerekvizitní
patch versus řešení uvnitř nového majoru; doba běhu migrace proti
odstávkovému oknu u největší instance.]}

% ========================================================================
\section{Antipatterny}
\label{sec:antipatterny}

\emph{[Stub. Squash bez zachování řetězu pro upgrade cestu. Ruční zásah
u zákazníka. On-prem s jedinou odchozí závislostí. Nulový počet hotfixů
jako cíl namísto výsledku.]}

% ========================================================================
\section*{Zdroje}
\addcontentsline{toc}{section}{Zdroje}

\emph{[Reference se zapojují přímo do textu u příslušných invariantů,
nikoli do samostatného appendixu. Nashromážděné:]}

\begin{itemize}[leftmargin=1.4em]
\item GitLab, charts/gitlab issue \#4522 --- vynucení poslední minor verze
      předchozí série jako vstupu do nového majoru. Nezávislý doklad
      k \theinvariant{} (viz sekce~\ref{sec:invarianty}).
\item GitLab, container-registry issue \#1516 --- rozbitý migrační stav
      při přeskočení post-deployment migrací.
\item Langfuse, self-hosting docs --- schéma a migrace verzované jako
      jeden celek; úprava schématu přesouvá odpovědnost na zákazníka.
\item Ardent Performance, \emph{Database Schema Migrations in 2026}
      (survey) --- srovnání přístupů u Sentry, Zulip a dalších.
\item Tenable, \emph{Software Release Lifecycle Policy} --- kadence
      a závazek k dostupnosti upgrade cest.
\item Microsoft, \emph{Finance + Operations (on-premises) version update
      policy} --- rozlišení kritických a nekritických aktualizací.
\item R. Mironov, \emph{Product Sprawl} --- popis symptomu z produktové
      perspektivy.
\end{itemize}

\end{document}
